\section{Experiments}
\label{sec:experiments}

\subsection{Experimental Setup}

\subsubsection{Dataset and Evaluation Metrics}
We evaluate our method on the Intentonomy dataset \cite{jia_intentonomy_2021}, a multi-label benchmark for visual intention recognition in social media scenarios. Following the standard split used in our experiments, the dataset contains 12,739 training images, 498 validation images, and 1,216 test images. Intentonomy additionally provides soft agreement annotations with values in $\{0, 0.33, 0.66, 1\}$ on the training set, which directly reflect the degree of annotator consensus and serve as the uncertainty signal $\omega_i$ in our framework.

We report Macro F1, Micro F1, Samples F1. For consistency with recent Intentonomy papers, the main comparison table additionally reports $\textit{Avg.}$ which is the average of three F1 scores. 

To better understand the effectiveness of our method, we further conduct experiments on subsets of the intent class. We follow the work of Jia et al. \cite{jia_intentonomy_2021} and break down intent classes according to content difficulty and dependency. Difficulty classes are measured by how much a standard ResNet-50 model outperforms the random results. Content classes are divided into object-dependent ({\itshape Object}-classes), context-dependent ({\itshape Context}-classes) and {\itshape Others} which depends on both foreground and background information. We emphasize the \textit{Hard} score throughout the analysis because it specifically measures performance on ambiguity-prone intent categories, where semantic overlap and subjective disagreement are most severe. Improvements on the hard subset therefore provide stronger evidence that a method is resolving the core difficulty of intention understanding rather than merely exploiting easy classes.

\subsubsection{Implementation Details}

Our default visual backbone is frozen CLIP ViT-L/14 \cite{radford_learning_2021}. The visual student is implemented as a lightweight residual MLP trained on top of a fixed SLR-C prior. For UTD, rationale text is generated offline by Qwen2.5-VL-7B-Instruct \cite{bai_qwen25-vl_2025} and then encoded by a frozen bge-large-en-v1.5 text encoder \cite{xiao_c-pack_2024}. The rationale-generation prompt template is provided in \ref{app:rationale_prompt}. The teacher consumes only the resulting text embeddings and is optimized with Asymmetric Loss (ASL) \cite{ridnik_asymmetric_2021} using $\gamma_{-}=2$, $\gamma_{+}=0$, and clipping value 0.05. The student is optimized with AdamW \cite{loshchilov_decoupled_2018}, learning rate $5\times10^{-4}$, weight decay $10^{-4}$, batch size 256, and dropout 0.1. The teacher hidden dimension is 1024 and the residual student hidden dimension is 768. The distillation temperature is set to 2.0. In SLR-C, scenario priors are constructed from Gemini 3.1 Pro-generated intent descriptions \cite{team_gemini_2025}, and local reranking is performed within Top-10 visual candidates. The scenario-prompt template is also provided in \ref{app:scenario_prompt}. During inference, class-wise thresholds are searched on the validation set and then fixed for test evaluation. Notably, UTD introduces \emph{zero extra inference cost}, since the text teacher is only used during training and is discarded at test time.


\subsection{Comparative Experiments}

\subsubsection{Comparison with state-of-the-art methods}
Table~\ref{tab:main_results} summarizes the main comparison on the testing set. More results on validation set can be found in \ref{app:validation_results}. The strong CLIP baseline already surpasses the previously reported best result by a clear margin, confirming that visual representation is no longer the primary bottleneck for this task. Once the backbone is sufficiently strong, the remaining challenge is how to map rich visual semantics into abstract and subjective intent labels. More importantly, \textit{UTD only} already improves over IntentMLM by a large margin while introducing zero extra inference cost, showing that robust training under subjective supervision is by itself sufficient to outperform recent multimodal large-model pipelines. Building on this, the complete SADIR framework further improves all four reported metrics on the testing split, reaching the strongest overall average score. This pattern is consistent with our motivation: SLR-C improves semantic discrimination at inference, whereas UTD stabilizes training under ambiguous supervision.

\begin{table}[H]
\centering
\small
\caption{Comparison with state-of-the-art methods. All numbers are percentages.}
\label{tab:main_results}
\begingroup
\setlength{\tabcolsep}{4pt}
\begin{tabular}{llcccc}
\hline
Method & Publication & Macro F1 & Micro F1 & Samples F1 & Avg. \\
\hline
Random & - & 6.94 & 7.18 & 7.10 & 7.07 \\
Visual \cite{he_deep_2016} & CVPR 2016 & 22.77 & 30.23 & 28.45 & 27.15 \\
\hline
\multicolumn{6}{l}{\textbf{Multi-label classification methods}} \\
HiMulConE \cite{zhang_use_2022} & CVPR 2022 & 28.43 & 41.50 & 42.15 & 37.36 \\
MultiGuard \cite{jia_multiguard_2022} & NeurIPS 2022 & 26.87 & 38.67 & 40.06 & 35.20 \\
SST \cite{chen_structured_2022} & AAAI 2022 & 27.33 & 42.21 & 42.62 & 37.39 \\
Query2Label \cite{liu_query2label_2021} & - & 32.12 & 44.64 & 45.15 & 40.64 \\
DualCoOp \cite{sun_dualcoop_2022} & NeurIPS 2022 & 31.06 & 44.42 & 40.40 & 38.63 \\
\hline
\multicolumn{6}{l}{\textbf{Label noise methods}} \\
Cleannet \cite{lee_cleannet_2018} & CVPR 2018 & 25.89 & 42.15 & 40.97 & 36.37 \\
LS \cite{lukasik_does_2020} & ICML 2020 & 24.32 & 42.77 & 43.27 & 36.79 \\
Dividemix \cite{li_dividemix_2019} & ICLR 2020 & 30.11 & 40.23 & 41.56 & 37.30 \\
\hline
\multicolumn{6}{l}{\textbf{Intent perception methods}} \\
Jia \cite{jia_intentonomy_2021} & CVPR 2021 & 23.98 & 31.28 & 31.39 & 28.88 \\
CPAD \cite{ye_cross-modality_2023} & TIP 2023 & 27.39 & 40.98 & 41.12 & 36.50 \\
HLEG \cite{shi_learnable_2023} & TAC 2023 & 32.77 & 44.69 & 45.54 & 41.00 \\
PIP-Net \cite{wang_prototype-based_2023} & TMM 2023 & 32.57 & 45.94 & 47.05 & 41.85 \\
LabCR \cite{shi_label-aware_2024} & TIP 2024 & 34.63 & 48.51 & 48.05 & 43.73 \\
IntCLIP \cite{leonardis_synergy_2025} & ECCV 2024 & 40.15 & 53.54 & 51.01 & 48.23 \\
% \hline
% \multicolumn{6}{l}{\textbf{Multimodal large model methods}} \\
% Qwen2-VL-7B & - & 21.85 & 27.06 & 33.48 & 27.46 \\
% Ovis1.6-Gemma2-9B & - & 36.14 & 34.52 & 33.93 & 34.86 \\
% GPT-4o & - & 39.48 & 40.57 & 41.45 & 40.50 \\
IntentMLM \cite{wang_unlocking_2026} & PR & 43.05 & 54.77 & 56.75 & 51.52 \\
\hline
CLIP ViT-L/14 \cite{radford_learning_2021} & - & 47.40 & 56.42 & 56.41 & 53.41 \\
SADIR (UTD only) & - & 50.31 & 57.97 & 59.07 & 55.78 \\
SADIR & - & \textbf{51.52} & \textbf{59.75} & \textbf{60.30} & \textbf{57.19} \\
\hline
\end{tabular}
\endgroup
\end{table}

\subsubsection{Performance on Different Difficulty Subsets}

To better understand the effectiveness of our method, we further evaluate performance on intent subsets with different difficulty levels. Following Jia et al. \cite{jia_intentonomy_2021}, we divide intent classes into \textit{Easy}, \textit{Medium}, and \textit{Hard} groups according to how much a standard ResNet-50 classifier outperforms random prediction. For previously published methods, we directly adopt the reported values from the difficulty-wise comparison figure provided in the original papers. For our models, the numbers are computed from the corresponding evaluation logs.

\begin{table}[t]
\centering
\caption{Difficulty-wise comparison on Intentonomy. All numbers are percentages.}
\label{tab:difficulty_results}
\begin{tabular}{lcccc}
\hline
Method & Easy & Medium & Hard & Avg. \\
\hline
Random & 19.86 & 7.11 & 2.81 & 9.93 \\
Visual \cite{he_deep_2016} & 54.64 & 24.92 & 10.71 & 30.09 \\
Jia (w/o hashtag) \cite{jia_intentonomy_2021} & 57.10 & 25.68 & 12.72 & 31.83 \\
Jia (hashtag) \cite{jia_intentonomy_2021} & 58.86 & 26.30 & 13.11 & 32.76 \\
Query2Label \cite{liu_query2label_2021} & 81.59 & 38.06 & 14.55 & 44.73 \\
CPAD \cite{ye_cross-modality_2023} & 69.47 & 31.50 & 8.54 & 36.50 \\
HLEG \cite{shi_learnable_2023} & 81.49 & 38.99 & 20.22 & 46.90 \\
PIP-Net \cite{wang_prototype-based_2023} & 70.73 & 33.63 & 14.33 & 39.56 \\
IntentMLM \cite{wang_unlocking_2026} & 76.26 & 46.87 & 27.39 & 50.17 \\
\hline
CLIP ViT-L/14 \cite{radford_learning_2021} & 80.99 & 53.32 & 28.45 & 54.25 \\
% Fixed SLR-C & 81.55 & 56.39 & 33.98 & 57.30 \\
% Ours: full model (standard KD) & 81.42 & 56.06 & \textbf{36.43} & \textbf{57.97} \\
SADIR & \textbf{82.53} & \textbf{56.60} & \textbf{34.60} & \textbf{57.91} \\
\hline
\end{tabular}
\end{table}

The difficulty-wise results provide two additional insights. First, our baseline is already substantially stronger than prior methods on all three difficulty groups, indicating that the new performance regime is not limited to a specific subset of classes. Second, SADIR further improves medium and hard groups most consistently, which is precisely where semantic ambiguity and supervision uncertainty are more pronounced.

\subsubsection{Performance of Different Content-Dependent Subsets}
Table~\ref{tab:content_results} reports results on content-dependent subsets. Following the protocol in prior work, intent classes are divided into \textit{Object-dependent}, \textit{Context-dependent}, and \textit{Others}. The first group mainly depends on salient foreground entities, the second relies more strongly on scene or environmental context, and the third requires combined reasoning over both foreground and background cues.

\begin{table}[t]
\centering
\caption{Comparison of results depending on different visual content. All numbers are percentages.}
\label{tab:content_results}
\begin{tabular}{lcccc}
\hline
Method & Object & Context & Others & Avg. \\
\hline
Random & 7.75 & 12.53 & 6.05 & 8.78 \\
Visual \cite{he_deep_2016} & 25.58 & 30.16 & 21.34 & 25.69 \\
Jia (w/o hashtag) \cite{jia_intentonomy_2021} & 28.15 & 28.62 & 22.60 & 26.46 \\
Jia (hashtag) \cite{jia_intentonomy_2021} & 29.66 & 32.48 & 22.61 & 28.25 \\
CPAD \cite{ye_cross-modality_2023} & 28.81 & 47.30 & 24.74 & 33.62 \\
PIP-Net \cite{wang_prototype-based_2023} & 39.20 & 39.53 & 27.59 & 35.44 \\
IntentMLM \cite{wang_unlocking_2026} & 50.13 & 49.69 & 39.75 & 46.52 \\
\hline
CLIP ViT-L/14 \cite{radford_learning_2021} & \textbf{61.80} & \textbf{61.69} & 40.59 & 54.69 \\
SADIR & 60.11 & 59.97 & \textbf{47.46} & \textbf{55.85} \\
\hline
\end{tabular}
\end{table}

The results again show that the strong CLIP baseline is already highly effective on object- and context-dependent intents. The major gain of SADIR appears in the \textit{Others} subset, where intent recognition requires more holistic reasoning over multiple weak cues rather than direct reliance on a single object or scene pattern.

\subsection{Ablation Studies}

\subsubsection{The Effectiveness of Rationale}

We first examine whether using rationale-style text is indeed more effective than generic image captioning. In an earlier shared-feature alignment setting, we compare caption supervision and rationale supervision under the same optimization protocol. As shown in Table~\ref{tab:caption_rationale_ablation}, the two are nearly indistinguishable, suggesting that simply replacing captions with free-form rationale text is not sufficient. This observation motivated the final UTD design, which uses structured rationale and agreement-aware distillation rather than shallow text-feature alignment alone.

\begin{table}[t]
\centering
\caption{Caption versus rationale in the earlier text-alignment setting. All numbers are percentages.}
\label{tab:caption_rationale_ablation}
\begin{tabular}{lcccc}
\hline
Text Source & Macro & Micro & Samples & Hard \\
\hline
Baseline & 48.23 & 52.04 & 52.94 & 29.71 \\
Caption & 48.34 & 47.55 & 47.54 & 30.28 \\
Rationale & 48.34 & 47.55 & 47.54 & 30.28 \\
\hline
\end{tabular}
\end{table}

We then ablate the internal structure of the rationale used by the final text teacher. Specifically, we compare three progressively richer versions: \textit{Step 1 only}, which contains observable visual evidence only; \textit{Step 1 with Step 2}, which additionally introduces contextual interpretation; and \textit{Full rationale}, which further adds counterfactual disambiguation cues. Table~\ref{tab:rationale_ablation} reports the results under dynamic gated distillation.

\begin{table}[t]
\centering
\caption{Ablation on the rationale construction strategy used by UTD. All numbers are percentages.}
\label{tab:rationale_ablation}
\begin{tabular}{cccccccc}
\hline
Step 1 & Step 2 & Step 3 & Macro & Micro & Samples & Avg. & Hard \\
\hline
$\checkmark$ &  &  & 49.81 & 56.39 & 57.65 & 54.62 & 29.28 \\
$\checkmark$ & $\checkmark$ &  & 49.69 & 54.45 & 55.97 & 53.37 & 28.82 \\
$\checkmark$ & $\checkmark$ & $\checkmark$ & \textbf{50.31} & \textbf{57.97} & \textbf{59.07} & \textbf{55.78} & \textbf{30.24} \\
\hline
\end{tabular}
\end{table}

The trend is revealing. Step 1 already provides useful visual facts, but adding Step 2 alone slightly degrades the overall classification metrics, suggesting that contextual interpretation without explicit disambiguation may introduce semantic drift. Once the counterfactual Step 3 is added, however, the rationale becomes significantly more effective. This supports our central claim that subjective intention recognition requires not only richer context, but also explicit boundary contraction among confusing alternatives.

\subsubsection{Effectiveness of Gated Distillation}

We next isolate the effect of UTD without using the SLR-C prior. Table~\ref{tab:distill_ablation} compares plain supervised learning, distillation without gating, and agreement-aware gated distillation.
\begin{table}[t]
\centering
\caption{Why gated distillation? All numbers are percentages.}
\label{tab:distill_ablation}
\begin{tabular}{ccccccc}
\hline
Backbone & Distillation & Gated & Macro & Samples & Avg. & Hard \\
\hline
$\checkmark$ & &  & 47.40 & 56.41 & 53.41 & 28.45 \\
$\checkmark$ &$\checkmark$ &  & 49.59 & 58.88 & 55.58 & \textbf{30.44} \\
$\checkmark$ &$\checkmark$ & $\checkmark$ & \textbf{50.31} & \textbf{59.07} & \textbf{55.78} & 30.24 \\
\hline
\end{tabular}
\end{table}

The results show that both distillation variants improve substantially over plain supervised learning, but agreement-aware gating provides the best overall balance in Macro and Samples F1. This supports our interpretation that dynamic gating acts as a form of \emph{on-demand distillation}: it suppresses over-regularization on reliable samples while reserving teacher guidance for genuinely ambiguous cases. The complementary evidence on weaker backbones is reported later in the robustness analysis.

\subsubsection{Modality Collapse Avoidance}

We further study whether the teacher should consume both image and text features, or only text features. Table~\ref{tab:teacher_modality_ablation} reports teacher-side oracle performance and student-side dynamic distillation performance as two separate method groups.
\begin{table}[t]
\centering
\caption{Teacher modality ablation. All numbers are percentages.}
\label{tab:teacher_modality_ablation}
\begin{tabular}{lccccc}
\hline
Method & Image & Text & Macro & Avg. & Hard \\
\hline
Oracle teacher &  & $\checkmark$ & 45.80 & 49.94 & \textbf{30.21} \\
Oracle teacher & $\checkmark$ & $\checkmark$ & 45.87 & 51.12 & 28.08 \\
\hline
Gated student &  & $\checkmark$ & \textbf{50.31} & \textbf{55.78} & \textbf{30.24} \\
Gated student & $\checkmark$ & $\checkmark$ & 48.57 & 53.93 & 28.58 \\
\hline
\end{tabular}
\end{table}

Interestingly, the multimodal oracle teacher is not weaker by itself, but it becomes a worse distillation target. This indicates that in a subjective recognition task, explicit image-text fusion inside the teacher may introduce noisy or redundant visual evidence, whereas a text-only semantic manifold provides cleaner dark knowledge for the student.

\subsubsection{Effectiveness of SLR-C and Residual Refinement}

We next study how SLR-C interacts with residual learning and distillation. As shown in Table~\ref{tab:slrc_ablation}, fixed SLR-C alone already provides a strong boost over the baseline, demonstrating that many remaining errors indeed lie in local candidate discrimination rather than visual representation learning. Training a residual student on top of this prior further improves the results, especially when semantic guidance from the text teacher is incorporated.
\begin{table}[t]
\centering
\caption{Component ablation of SADIR. ``Res.'' denotes the residual student, and ``Dyn.'' denotes agreement-aware dynamic gating. All numbers are percentages.}
\label{tab:slrc_ablation}
\begin{tabular}{ccccccccc}
\hline
SLR-C & Student & Distillation & Gated & Macro & Micro & Samples & Avg. & Hard \\
\hline
$\checkmark$ & &  &  & 51.08 & 58.81 & 58.02 & 55.97 & 33.98 \\
$\checkmark$ &$\checkmark$ &  &  & 50.41 & 57.70 & 57.56 & 55.22 & 34.77 \\
$\checkmark$ &$\checkmark$ & $\checkmark$ &  & \textbf{51.77} & 58.97 & 58.99 & 56.58 & \textbf{36.43} \\
$\checkmark$ &$\checkmark$ & $\checkmark$ & $\checkmark$ & 51.52 & \textbf{59.75} & \textbf{60.30} & \textbf{57.19} & 34.60 \\
\hline
\end{tabular}
\end{table}

Two observations are noteworthy. First, simply adding a residual branch with hard supervision is not sufficient to guarantee gains on all metrics, indicating that the fixed prior already captures much of the easy signal. Second, once semantic distillation is introduced, the residual student becomes consistently beneficial, showing that UTD and SLR-C address different failure modes and are therefore naturally complementary.

\subsubsection{Heterogeneous Priors in SLR-C}
Finally, we compare different textual priors used by SLR-C. We consider three semantic sources: \textit{Lexical}, which uses short class-name phrases; \textit{Canonical}, which uses the official intent descriptions provided by Intentonomy; and \textit{Scenario}, which uses Gemini-generated visual scenarios. For the ensemble setting, we average the normalized logits from the selected sources before local reranking. Table~\ref{tab:prior_ablation} shows that \textit{Scenario} prior is the strongest single source, while the heterogeneous lexical-canonical combination provides the best non-scenario alternative.
\begin{table}[t]
\centering
\caption{Ablation on heterogeneous textual priors for SLR-C. All numbers are percentages.}
\label{tab:prior_ablation}
\begin{tabular}{cccccccc}
\hline
Lexical & Canonical & Scenario & Macro & Micro & Samples & Avg. & Hard \\
\hline
$\checkmark$ &  &  & 47.06 & -- & -- & -- & 30.53 \\
 & $\checkmark$ &  & 47.76 & -- & -- & -- & 31.01 \\
 &  & $\checkmark$ & 51.28 & \textbf{59.13} & 58.47 & \textbf{56.29} & 33.98 \\
$\checkmark$ & $\checkmark$ &  & 50.65 & 57.37 & 57.46 & 55.16 & 32.38 \\
$\checkmark$ & $\checkmark$ & $\checkmark$ & \textbf{51.49} & 58.76 & \textbf{58.54} & 56.26 & \textbf{35.08} \\
\hline
\end{tabular}
\end{table}

This result is consistent with our semantic interpretation of the task. Lexical prior is compact but under-expressive, canonical prior is more descriptive but still abstract, while scenario prior translates latent intention into typical visual situations and therefore best bridges the gap between visual content and subjective semantics. The lexical-canonical ensemble remains competitive, showing that heterogeneous priors are beneficial when richer scenario descriptions are not available. When all three sources are combined, the resulting ensemble improves substantially over lexical-canonical priors and further raises the hard score to 35.08, but its overall result remains slightly below the scenario-only prior. This suggests that scenario descriptions already capture the most decisive semantics, while additional sources mainly provide complementary robustness on difficult classes.

\subsection{Robustness and Generalization}

\subsubsection{Backbone Independence}

To examine whether SADIR depends entirely on a strong backbone, we also test the full pipeline with a weaker ResNet101 student. The results are shown in Table~\ref{tab:backbone_ablation}. The full method remains functional, but the gains become much smaller, and fixed SLR-C becomes the dominant source of performance. This suggests that the framework is portable to weaker backbones, yet its full benefit is realized only when the visual backbone is already reasonably strong.

\begin{table}[t]
\centering
\caption{Backbone sensitivity comparing baseline and full model. All numbers are percentages.}
\label{tab:backbone_ablation}
\begin{tabular}{lcccccc}
\hline
Backbone & Method & Macro & Micro & Samples & Avg. & Hard \\
\hline
ResNet101 & Baseline & 36.56 & 41.87 & 43.10 & 40.51 & 18.99 \\
ResNet101 & Full model & 41.53 & 51.61 & 51.80 & 48.31 & 22.10 \\
\hline
CLIP ViT-L/14 & Baseline & 47.40 & 56.42 & 56.41 & 53.41 & 28.45 \\
CLIP ViT-L/14 & Full model & \textbf{51.52} & \textbf{59.75} & \textbf{60.30} & \textbf{57.19} & \textbf{34.60} \\
\hline
\end{tabular}
\end{table}

\subsection{Visualization and Case Study}
